% CLAUDE WR REPORT::WR-BRIDGE-PINCER-0001::2026-11-15
\documentclass[11pt]{article}
\usepackage{amsmath,amssymb,amsthm}
\usepackage[margin=1in]{geometry}
\usepackage{hyperref}
\newtheorem{theorem}{Theorem}
\newtheorem{lemma}[theorem]{Lemma}
\newtheorem{corollary}[theorem]{Corollary}
\theoremstyle{definition}
\newtheorem{definition}[theorem]{Definition}
\newcommand{\Wbar}{\overline W}
\title{The B3 Bridge for Worst-Case Minimum Derivation Cost:\\
A Dichotomy, a Monotone Pincer, and the $\exists/\forall$ Collapse}
\author{Claude scheduled routine (worst-case derivation directive)}
\date{2026-11-15}
\begin{document}
\maketitle

\section*{Status}
\textbf{P versus NP remains OPEN.} Lemmas 1--3 and Corollary 4 are
\textbf{PROVED/EXACT} and elementary; Theorem 5 is a \textbf{KNOWN THEOREM}
(Krajíček--Pudlák); Corollary 6 is \textbf{PROVED/EXACT} given Theorem~5.
No novelty is claimed beyond packaging: each piece is folklore or follows in a
few lines from standard results. The point of the note is to show exactly how
much the ``weakest sufficient bridge''
\[
\textbf{(B3)}\qquad \mathrm{SAT}\in\mathrm P\ \Longrightarrow\ \Wbar_R(N)\le\mathrm{poly}(N)
\]
can buy, for each choice of the fixed system $R$. No computation was run;
the results are purely mathematical.

\section*{Model and quantifiers}
$R$ is a Cook--Reckhow refutation system for CNF: a polynomial-time verifier
$V_R(F,\pi)$, sound and complete, fixed before any instance or size is chosen.
\[
W_R(F)=\min_{\pi:\,V_R(F,\pi)=1}\mathrm{cost}_R(\pi),\qquad
\Wbar_R(N)=\max_{F\in\mathrm{UNSAT},\ |\langle F\rangle|\le N}W_R(F).
\]
\textbf{Weight assumption (W).} $\mathrm{cost}_R(\pi)$ and the encoded length
$|\pi|$ are polynomially related in both directions. This holds for line or
bit cost in Resolution, CP with binary coefficients, PC, Frege and EF. It
\emph{fails} for unit-cost DAG nodes whose lines can be arbitrarily large
(for example, unbounded GF(2) rows or circuits counted as one node). For those
measures, Lemma~3 needs a separate size bridge.
Write $\mathrm{LB}(R)$ for ``$\Wbar_R$ is superpolynomial'' and $\mathrm{PB}(R)$
for ``$\Wbar_R\le\mathrm{poly}$'' ($R$ is p-bounded).
The inner minimum makes the best derivation free; search cost never enters.

\section{Results}

\begin{lemma}[B3 dichotomy]\label{l:dich}
For every fixed $R$, $\mathrm{B3}(R)\iff(\mathrm P\ne\mathrm{NP})\vee\mathrm{PB}(R)$.
\end{lemma}
\begin{proof}
This is the material implication $(\mathrm P=\mathrm{NP})\Rightarrow\mathrm{PB}(R)$.
\end{proof}

\begin{corollary}[Known-lower-bound systems give no leverage]\label{c:known}
If $\mathrm{LB}(R)$ is already proved, then $\mathrm{B3}(R)\iff\mathrm P\ne\mathrm{NP}$.
This applies to Resolution (Haken: PHP; Urquhart: Tseitin), Cutting Planes
(Pudlák: clique--colouring), Polynomial Calculus (Razborov; Ben-Sasson--Impagliazzo),
bounded-depth Frege (Ajtai; Beame--Impagliazzo--Pitassi; Krajíček--Pudlák--Woods)
and depth-restricted or regular Res$(\oplus)$.
For each of these systems, proving B3 \emph{is} proving P$\ne$NP. The bridge is no
weaker than the goal.
\end{corollary}

\begin{lemma}[Monotone pincer]\label{l:mono}
If $R'$ p-simulates $R$, then $\mathrm{LB}(R')\Rightarrow\mathrm{LB}(R)$ and
$\mathrm{B3}(R)\Rightarrow\mathrm{B3}(R')$.
\end{lemma}
\begin{proof}
A p-simulation maps each $R$-refutation to an $R'$-refutation of polynomially
related size. So $\mathrm{PB}(R)\Rightarrow\mathrm{PB}(R')$, and both claims
follow from Lemma~\ref{l:dich}.
\end{proof}
So lower bounds get easier as $R$ gets weaker, and B3 gets easier as $R$ gets
stronger. The route ``prove $\mathrm{LB}(R)$, then prove $\mathrm{B3}(R)$
without proving P$\ne$NP first'' needs one $R$ whose lower bound is currently
\emph{unknown}. By Corollary~\ref{c:known}, every system with a known lower bound
falls outside that region.

\begin{lemma}[Under P=NP, p-bounded $\iff$ p-optimal]\label{l:opt}
Assume (W) and $\mathrm P=\mathrm{NP}$. Then $R$ is p-bounded iff $R$ is
p-optimal (it p-simulates every refutation system with a polynomial-time
translation). Hence $\mathrm{B3}(R)\iff[\mathrm P=\mathrm{NP}\Rightarrow R\text{ p-optimal}]$.
\end{lemma}
\begin{proof}
($\Leftarrow$) Take a polynomial-time SAT decider $A$. The system $Q_A$ accepts
$(F,1^{|F|})$ iff $A$ rejects $F$. It is a sound, complete Cook--Reckhow system
and it is p-bounded. A p-optimal $R$ p-simulates $Q_A$, so $R$ is p-bounded.
($\Rightarrow$) Suppose $W_R(F)\le p(|F|)$. The language
$\{(F,\sigma):\exists\pi\supseteq\sigma,\ |\pi|\le p(|F|),\ V_R(F,\pi)=1\}$ is
in NP, hence in P. Prefix self-reduction builds an $R$-refutation of any UNSAT
$F$ in polynomial time, whatever foreign proof is given as input. That is a
polynomial-time translation, so $R$ is p-optimal.
\end{proof}

\begin{theorem}[KNOWN THEOREM: Krajíček--Pudlák 1989/90; cf.\ Cook 1975]\label{t:kp}
For every refutation system $Q$, EF p-simulates $Q$ iff EF has polynomial-size
proofs of the reflection tautologies $\|\mathrm{Ref}_Q\|_n$.
\end{theorem}

\begin{corollary}[$\exists/\forall$ collapse for EF]\label{c:ea}
The following are equivalent:
(i) $\mathrm{B3}(\mathrm{EF})$;
(ii) $\mathrm P=\mathrm{NP}\Rightarrow$ for \emph{some} polynomial-time correct
SAT decider $A$, EF has poly-size proofs of $\|\mathrm{Ref}_{Q_A}\|_n$;
(iii) $\mathrm P=\mathrm{NP}\Rightarrow$ for \emph{every} polynomial-time correct
SAT decider $A$, EF has poly-size proofs of $\|\mathrm{Ref}_{Q_A}\|_n$.
\end{corollary}
\begin{proof}
(iii)$\Rightarrow$(ii) is trivial under $\mathrm P=\mathrm{NP}$, since some $A$ exists.
(ii)$\Rightarrow$(i): Theorem~\ref{t:kp} gives EF $\ge Q_A$, and $Q_A$ is p-bounded.
(i)$\Rightarrow$(iii): by Lemma~\ref{l:opt}, EF is p-optimal, so it simulates every
$Q_A$; then apply Theorem~\ref{t:kp}. The same argument works for any system
that satisfies the Krajíček--Pudlák reflection characterisation, e.g.\
EF$+\|\mathrm{Ref}_T\|$ for a fixed theory $T$.
\end{proof}
\textbf{Reading.} For a strong $R$, B3 is exactly the corpus's OPEN
``affordable reflection'' obligation (codex 2026-11-11\_1331). It needs short
$R$-proofs of the soundness of the \emph{hypothetical} fast algorithm's UNSAT
answers. The corollary shows that the choice of algorithm does not matter: if
one fast algorithm has short soundness proofs, every fast algorithm does. So a
B3 proof cannot exploit the structure of a particular cleverly chosen $A$. It
has to work for an arbitrary, possibly obfuscated $A$. A bare execution trace
of $A$ is never enough without reflection (Theorem~\ref{t:kp}), which matches
the earlier LSPV-5 REFEREE294 countercontrol.

\section{Directive gates}
\begin{itemize}
\item \textbf{Free-strategy clause.} Respected: $W_R$ is a minimum over proofs.
Lemma~\ref{l:opt} uses search only under the hypothesis P=NP, and only to build
a translation; no lower-bound argument charges for search.
\item \textbf{DISJ sanity.} Not applicable: no communication or rank invariant is claimed.
\item \textbf{Relativization.} Lemmas 1--3 relativize. They are purely logical
or use only self-reducibility. They therefore cannot by themselves separate
anything, which is consistent with their status as reformulations.
\item \textbf{Simultaneous DTISP / configuration counting.} Not used. No
machine time or space bound is inferred from proof size.
\item \textbf{World ledger.} Lemma~\ref{l:dich} with LB known is a statement
\emph{equivalent to P$\ne$NP}, i.e.\ a reformulation and not progress
(root AGENTS.md world-ledger guidance).
\end{itemize}

\section{Hostile controls and falsifiers}
\begin{enumerate}
\item \emph{Tseitin on expanders.} Resolution needs exponential size, yet
Gaussian elimination decides them in polynomial time. This refutes pointwise
B1 for Resolution. It does \emph{not} refute B3 for Resolution, which by
Corollary~\ref{c:known} is equivalent to P$\ne$NP. This is consistent with
LSPV-1/334--335.
\item \emph{Falsifier F1.} Any claimed proof of $\mathrm{B3}(R)$ for a system
listed in Corollary~\ref{c:known} that does not already prove P$\ne$NP is wrong.
\item \emph{Falsifier F2.} Any claimed proof of $\mathrm{B3}(\mathrm{EF})$ that
works for one specially structured $A$ but not for arbitrary $A$ is wrong by
Corollary~\ref{c:ea}. For example, it might rely on $A$ being a resolution or
DPLL procedure.
\item \emph{Falsifier F3 (weights).} Under unit node cost with unbounded line
size (as in the corpus's Rmix341 GF(2) rows), assumption (W) can fail, and
Lemma~\ref{l:opt} must be re-proved for the actual cost measure. Rmix341's
linear optimum $4h-1$ on Tseitin is an example of a small node count with
$\Theta(h)$-bit rows.
\end{enumerate}

\section{Literature check (this run)}
\begin{itemize}
\item Unrestricted Res$(\oplus)$: superpolynomial size lower bounds are still
not known (Efremenko--Itsykson, Nov.\ 2025, per ECCC/arXiv). Depth-restricted
bounds keep improving: Bhattacharya--Chattopadhyay arXiv:2507.23008 (depth
$N^{1+\epsilon}$); Byramji--Impagliazzo arXiv:2511.20023 (depth $N^{2-\epsilon}$,
size $\exp(\Omega(N^\epsilon))$); ECCC TR25-118 (size $\exp(\Omega(n^3/D^2))$);
CCC 2026 (depth $o(n^2/\log^4 n)$).
\emph{Consequence for the pincer:} Res$(\oplus)$, and the corpus's affine mixed
calculus, which it plausibly p-simulates, is one of the weakest natural systems
whose LB is open. It is therefore the most natural current point of the region
in Lemma~\ref{l:mono}.
\item Pich--Santhanam (ECCC 2023/199): a conditional EF-to-P$\ne$NP transfer that
uses extra hypotheses. It is adopted as the published form of the OPEN reflection target.
\item ECCC TR26-109 (2026): reflection principles and implicit proof systems. It
records that there is still no concrete non-reflection tautology family believed
hard for EF but easy for a stronger system.
\item Cook--Reckhow 1979 (a p-bounded system exists iff NP=coNP); Krajíček--Pudlák 1989/90.
\end{itemize}

\section*{CORPUS REVIEWED}
AGENTS.md (including the 2026-11-15 SAT-compression directive);
research/coordination/2026-11-02\_1204\_WORST\_CASE\_MINIMUM\_DERIVATION\_FIVE\_ROLE\_REFOCUS.md
and LSPV-2/341 inside it; LSPV-1/335 (worst-case core and explicit-generator barrier);
codex 2026-11-11\_1203 (SAT polynomial assumption audit) and 2026-11-11\_1331
(transfer recovery and reflection); a repository search for p-optimal and reflection material.

\section*{DUPLICATION / CONTRADICTION AUDIT}
Lemma~\ref{l:dich} and Theorem~\ref{t:kp} are implicit in LSPV-1/335 and in the
codex reflection sections. The p-bounded$\iff$p-optimal equivalence under P=NP and
the $\exists/\forall$ collapse (Corollary~\ref{c:ea}) were not found stated
explicitly in the corpus searched. Both are folklore-level, and no novelty is
claimed. No contradiction with the corpus was found.

\section*{WHAT CHANGED RELATIVE TO ROOT MAIN}
Nothing in root main.tex was edited. This routine owns no LSPV/TMR ownership
region, so under the ownership rules it does not edit main.tex.

\section*{ADOPTED / REJECTED / RETIRED PRIOR COMPONENTS}
Adopted: the B1/B2/B3 hierarchy, the explicit-generator barrier (LSPV-1/335) and
affordable reflection as OPEN. Retired: none. Narrowed: B3 for any system with a
known lower bound is classified as equivalent to P$\ne$NP (Corollary~\ref{c:known}).

\section*{STALE OR OVERSTATED CLAIMS FOUND}
None in the files read.

\section*{Next missing implication}
Choose $R\in\{\mathrm{Res}(\oplus),\ \text{Rmix341 with bit-weighted rows}\}$. Then
either (a) prove LB($R$) for a family that is \emph{not} explicitly generated in
polynomial time (the LSPV-1/335 barrier forbids an explicit family if B1 is also
wanted), or (b) show that EF-reflection for arbitrary fast $A$ is unachievable
in $R$, which would close the B3 route for $R$. Both remain OPEN.

\end{document}
% END_OF_VERIFIED_REPORT::CLAUDE-WR::WR-BRIDGE-PINCER-0001
